\documentclass[11pt]{amsart}
\usepackage[T1]{fontenc}
\usepackage[utf8]{inputenc}
\usepackage{lmodern}
\usepackage{amsmath,amssymb,amsthm,mathtools}
\usepackage{enumitem}
\usepackage[hidelinks]{hyperref}

\newtheorem{theorem}{Theorem}[section]
\newtheorem{lemma}[theorem]{Lemma}
\newtheorem{proposition}[theorem]{Proposition}
\newtheorem{remark}[theorem]{Remark}
\theoremstyle{definition}
\newtheorem{definition}[theorem]{Definition}

\begin{document}

\title[Switch Corridors: Incomplete Architecture]{Switch Corridors and Boundary-Rigid Lenses in Even Dumbbell Thrackles: An Incomplete Proof Architecture and Gap Audit}
\author{Luca Blanchi}
\date{June 2026}

\begin{abstract}
This note does not prove Conway's thrackle conjecture. It records an
AI-assisted proof architecture that was initially drafted as a possible proof,
together with a subsequent audit of the places where the argument is sound,
conditional, or presently incomplete.

The program tries to reduce the internal even-dumbbell problem to a
combination of switch-corridor regularization, support-gate disk
classification, rank-one analysis at the common vertex, and certified
edge-removal reductions. The audit shows that several bookkeeping and local
geometric components can be made precise, but the central support-gate router
is not closed. In particular, the remaining proof would still need a complete
classification of label-essential support-gate disks containing original
vertices, or an equivalent theorem producing certified edge-removal
configurations with controlled triangles.

The useful output of the note is therefore not a finished proof, but a
structured failure audit: it identifies which reductions are legitimate, which
old formulations were too strong, and which residual lemmas would be needed
before the strategy could become a proof.
\end{abstract}

\maketitle

\section*{Current Status}

This manuscript is a preliminary AI-assisted research note. It should be read
as an incomplete proof architecture, not as a proof of Conway's thrackle
conjecture and not as a proof that irreducible even dumbbell thrackles are
$T_3$.

The strongest honest statement at the present stage is:

\[
\text{the strategy has a coherent blueprint, but its central router is open.}
\]

The main unresolved point is the classification of label-essential
support-gate disks containing original vertices. The current audit disables
all automatic uses of a smaller-core or R2-type reduction unless an explicit
external certificate is supplied. The viable certified reduction route is R3,
but R3 requires both a degree-two three-path candidate and an empty or
controlled edge-removal triangle. Those certificates are not yet available in
all residual cases.

\section*{Declaration of Generative AI and AI-Assisted Technologies}

During the preparation of this work, ChatGPT, by OpenAI, was used to assist
with mathematical drafting, formalization, review, editing, and gap auditing.
This work is shared as a preliminary AI-assisted mathematical note. The
mathematical content may have been only partially reviewed and may contain
errors; it should not be treated as peer-reviewed or as a fully verified
manuscript.

\section{Purpose of the Note}

The original draft attempted to prove Conway's thrackle conjecture by the
following high-level route:

\[
\text{Conway false}
\Rightarrow
\text{even dumbbell thrackle}
\Rightarrow
\text{irreducible even dumbbell thrackle}
\Rightarrow
T_3
\Rightarrow
|E|\le |V|.
\]

The external inputs in that route are standard structural ingredients: the
even-dumbbell reduction and the theorem that $T_3$-thrackles satisfy Conway's
inequality. The internal target was:

\[
\text{every irreducible even dumbbell thrackle is } T_3.
\]

This note no longer claims that internal target. Instead, it records the
present state of the attempted internal proof.

\section{Original Blueprint}

Let

\[
F_{2a,2b}=C_{2a}\vee C_{2b}
\]

be an even dumbbell, with cycles $A=C_{2a}$ and $B=C_{2b}$ sharing the vertex
$v$.

The intended strategy was:

\begin{enumerate}
\item Use a domain-incidence argument to select ordinary vertices $p,q$ on
      one cycle with disjoint incident-domain pairs.
\item Study the two complementary $p$-to-$q$ paths on that cycle.
\item Extract switch-corridor or support-gate witnesses carrying the
      nontrivial labelled transport.
\item Regularize clean $v$-free parts into support-switch ladders or closed
      support-switch trains.
\item Treat the path through $v$ by a singular rank-one analysis rather than
      as an ordinary support-switch cell.
\item Classify the label-essential support-gate loops that arise from raw
      closures, repeated cores, train normalization, or rank-one transfer.
\item Conclude by contradiction from either an impossible train, a legitimate
      complexity-lowering reduction, or a certified edge-removal operation.
\end{enumerate}

After audit, the blueprint remains useful, but the last two items are not
complete.

\section{Corrections Introduced by the Audit}

\subsection{No automatic proof of the internal theorem}

The old draft repeatedly stated that every irreducible even dumbbell thrackle
is $T_3$. That statement is now only the desired target. It is not established
by the current argument.

Any final section deriving Conway's conjecture from the internal theorem must
therefore be read conditionally:

\[
\text{if the internal theorem were proved, then the external reduction would close Conway.}
\]

\subsection{Anchored core carriers replace arbitrary minimal corridors}

The old ``minimal corridor'' language was too loose. A globally minimal
nontrivial disk corridor may have endpoints in the middle of a raw trace and
need not be anchored at the selected vertices $p,q$. That is a problem because
the ladder-closing step requires compatible endpoint data at $p$ and $q$.

The corrected object is an anchored core carrier: a witness that preserves the
actual $p$-to-$q$ core trace, or else explicitly routes failures to a smaller
closed walk, support-gate loop, R1 splice, rank-one marker case, train, or
typed normality pathology.

\subsection{The common vertex is singular}

The common vertex $v$ cannot be treated as an ordinary support-switch cell.
After auditing the local rotation, the separated rank-zero case is ruled out
by a parity argument, but the surviving local model is still rank-one, not
full-rank.

With local sectors denoted by

\[
S_A,\quad S_1,\quad S_B,\quad S_2,
\]

the ordinary continuations are

\[
c_A:S_A\to S_A,\qquad c_B:S_B\to S_B,
\]

while the mixed transitions are marked singular transitions

\[
m_A:S_2\Rightarrow S_1,\qquad m_B:S_1\Rightarrow S_2.
\]

The arrows $\Rightarrow$ are not ordinary lane continuations. They record
rank-one marker uses at $v$.

\subsection{The M-distinct boundary-component formulation was wrong}

The old M-distinct argument used a regular-neighborhood boundary component of
the planarized drawing $A\cup B$. That formulation is not reliable: planarized
crossings change the first Betti number, and distinct mixed sectors need not
lie on a single geometric face-boundary component.

The corrected object is a finite closed mixed support-gate state trace

\[
L_{\mathrm{mix}}=A_{\mathrm{mix}}\cup B_{\mathrm{mix}},
\]

where $A_{\mathrm{mix}}$ and $B_{\mathrm{mix}}$ carry the two opposite marked
singular transitions. Cancellation then requires a marked-collapse or
same-relation transfer argument; it is not a local ordinary cancellation.

\subsection{Labelled witnesses must be typed}

The reduction language now uses typed labelled witnesses. A boundary relation
is not just an order relation between endpoints; it must remember typed
terminal states, global domain labels, orientation data, and marker records.

In particular, two appearances on the same geometric support or gate interval
are not automatically the same state. They are the same state only if the
active support/gate, side, domain-label, orientation, boundary-role, and marker
data agree.

This prevents false collapses, false gate identifications, and false
repeated-support reductions.

\subsection{Pathology is not a terminal proof word}

The word ``pathology'' is now only shorthand for a finite routed ledger:

\begin{itemize}
\item terminal lane;
\item returning lane;
\item branching or merging;
\item gate identification;
\item repeated support;
\item double crossing;
\item cell overlap.
\end{itemize}

Each such event must immediately route to a labelled collapse, a
same-relation splice, a smaller label-essential support-gate loop, a smaller
train, a rank-one marker case, a certified R3 configuration, or a direct
thrackle contradiction. If no such route is named, the case is unresolved.

\subsection{R2 is disabled without an external certificate}

The old phrase ``relevant inner domain'' hid too much. In a pure even
dumbbell, a disk containing original vertices does not automatically contain a
proper smaller even dumbbell obstruction.

The corrected policy is:

\[
\text{do not invoke R2 unless an explicit compression or replacement-edge certificate is supplied.}
\]

Under the present pure-dumbbell blueprint, R2 is a disabled placeholder.
Apparent R2b edge-removal cases should be treated as R3 candidates. Apparent
R2c replacement-edge cases require a separate theorem.

\subsection{R3 needs a controlled triangle}

The certified R3 route is valid only when an exact edge-removal configuration
is exhibited. It requires:

\begin{enumerate}
\item a degree-two three-path $x_0x_1x_2x_3$;
\item the correct crossing of the first and third edges;
\item an empty, inessential, or otherwise controlled triangle bounded by the
      middle edge and the corresponding subarcs;
\item persistence of the relevant non-$T_3$ obstruction after edge removal.
\end{enumerate}

Long path-like components avoiding $v$ and long arms in $v$-containing disks
give plausible degree-two three-path candidates. They do not by themselves
give the controlled triangle certificate.

\section{Blocks That Look Blindable}

The following blocks appear solid enough to be incorporated into a future
clean manuscript, subject to ordinary checking and notation cleanup.

\subsection{Rotation parity at the common vertex}

The rank-zero local rotation at $v$ is ruled out in even dumbbells by an
immersed-arc parity argument. After deleting a small disk around $v$, the two
cycles give arcs with an even number of mutual crossings, so the endpoint
pattern on the boundary cannot alternate in the separated way.

\subsection{Typed rank-one marker at $v$}

The local interface at $v$ is correctly modelled by ordinary continuations
$c_A,c_B$ and marked singular transitions $m_A,m_B$. The mixed transitions
must remain marked until a separate cancellation or same-relation replacement
is certified.

\subsection{Labelled witness operations}

The abstract minimality operations are sound once their hypotheses are met:

\begin{itemize}
\item full-labelled identity deletion;
\item same-relation two-terminal splice;
\item rank-defect-lowering replacement.
\end{itemize}

The important correction is that these operations preserve full typed labelled
relations, not merely endpoint order.

\subsection{Cell-only support-gate disks}

For a cell-clean support-gate disk containing only support-switch cells, the
finite incidence graph argument is essentially graph-theoretic. Under typed
normality and labelled-collapse hypotheses, it routes cyclic components to
smaller trains and acyclic components to smaller loops, collapses, or typed
normality failures.

\subsection{Pure dumbbell port graph bookkeeping}

Inside a pure even dumbbell, original-vertex disk contributions have no hidden
automatic smaller-core form. They route to path pieces, $v$-arms,
support-gate boundary closures, whole-cycle or whole-dumbbell cases, certified
R3, explicit external R2, or the typed normality ledger.

\section{Central Open Block}

The central support-gate theorem should currently be stated only as a
conditional router:

\[
\begin{aligned}
\text{innermost label-essential support-gate loop}
\Rightarrow {}&
\text{smaller train}\\
&\text{or typed ledger pathology}\\
&\text{or R1 splice}\\
&\text{or rank-one }v\text{ marker package}\\
&\text{or labelled collapse}\\
&\text{or certified R3}\\
&\text{or explicit external R2 obligation.}
\end{aligned}
\]

This is not yet a proof, because two outputs are not automatically
contradictions:

\begin{enumerate}
\item R2 is disabled unless a new external certificate is supplied.
\item R3 is certified only after both the degree-two three-path and the
      controlled triangle have been produced.
\end{enumerate}

Thus the main missing theorem is one of the following:

\begin{enumerate}
\item prove that every residual original-vertex support-gate disk supplies a
      certified R3 edge-removal triangle;
\item prove a new legitimate external R2 compression or replacement-edge
      theorem;
\item reroute all residual cases to R1, rank-one marker descent, labelled
      collapse, smaller train, or typed normality contradictions.
\end{enumerate}

\section{What the Complexity Ideas Contribute}

The complexity viewpoint remains useful, but it is not currently a complete
engine for the proof. Its useful contribution is proof accounting:

\begin{itemize}
\item every collapse or splice must preserve the full typed labelled boundary
      relation;
\item every replacement must lower a declared lexicographic complexity;
\item rank-one marker defects at $v$ must be counted before ordinary
      support-gate length;
\item marker-only normal $v$-ears must not be fed back into the ordinary
      support-gate theorem as if they were new exterior essentiality;
\item a ``smaller'' object must be smaller in the appropriate global,
      witness, loop, train, or marker hierarchy.
\end{itemize}

This accounting prevents several false reductions. It has not yet produced
the missing geometric certificate needed to close the proof.

\section{Current Verdict}

The present document should be classified as:

\[
\text{incomplete proof architecture plus gap audit.}
\]

It contains useful components and a more honest map of the proof landscape,
but it does not prove Conway's thrackle conjecture. In particular, it does
not prove that every irreducible even dumbbell thrackle is $T_3$.

The next mathematical target is the residual support-gate disk problem:
either certify R3 in all remaining cases, or replace the R3 branch by a
different legitimate descent theorem.

\begin{thebibliography}{99}

\bibitem{LPS97}
L. Lovasz, J. Pach, and M. Szegedy,
\emph{On Conway's thrackle conjecture},
Discrete \& Computational Geometry 18(4) (1997), 369--376.

\bibitem{CN00}
G. Cairns and Y. Nikolayevsky,
\emph{Bounds for generalized thrackles},
Discrete \& Computational Geometry 23(2) (2000), 191--206.

\bibitem{MN18}
G. Misereh and Y. Nikolayevsky,
\emph{Annular and pants thrackles},
Discrete Mathematics \& Theoretical Computer Science 20(1) (2018), Article \#16.

\bibitem{FP11}
R. Fulek and J. Pach,
\emph{A computational approach to Conway's thrackle conjecture},
Computational Geometry: Theory and Applications 44(6--7) (2011), 345--355.

\bibitem{Xu21}
Y. Xu,
\emph{A new upper bound for Conway's thrackles},
Applied Mathematics and Computation 389 (2021), Article 125573.

\end{thebibliography}

\end{document}
